Take a square piece of paper of thickness $d$ and side lengths $l=1$ with barycenter at the origin. First, fold the paper randomly and endow the two resulting sides with different topologies. Fold it back to the original structure. Now fold it three times along the longest edge, i.e., half the longest edge. All folding is performed in Euclidean space $\mathbb R^3$.

The structure you have now will have layers on top of each other endowed with different topologies. Now take one of the vertices and fold it to the middle of the longest edge.

Now subdivide the resulting structure into a triangle and a rectangular shape. Puncture the triangle at its barycenter with a pointy steel rod. Unfold the structure again to its original shape with side length 1.

Group the holes according to their distance to the origin. Remember that you have two topologies. This gives the tuple of the cardinality of the distance classes.

Now take the length of this tuple and define it as $c$. Organize all holes on the paper, which has thickness $d$, in a cyclic graph. Let $c$ be the number of possible colors for the edges.
Compute the number of monochromatic perfect matchings of the graph.
